%% Template for ENG 401 reports
%% by Robin Turner
%% Adapted from the IEEE peer review template

%
% note that the "draftcls" or "draftclsnofoot", not "draft", option
% should be used if it is desired that the figures are to be displayed in
% draft mode.

\documentclass[peerreview]{IEEEtran}
\usepackage{cite} % Tidies up citation numbers.
\usepackage{url} % Provides better formatting of URLs.
\usepackage[utf8]{inputenc} % Allows Turkish characters.
\usepackage{booktabs} % Allows the use of \toprule, \midrule and \bottomrule in tables for horizontal lines
\usepackage{graphicx}


\hyphenation{op-tical net-works semi-conduc-tor} % Corrects some bad hyphenation 



\begin{document}
%\begin{titlepage}
% paper title
% can use linebreaks \\ within to get better formatting as desired
\title{Review: Coffee Ingredient Usage Forecasting Using
MySQL, MongoDB, and Moving Average}


% author names and affiliations
\date{13.07.2026}

% make the title area
\maketitle

\section{Summary Rating}
4 - accept 

\section{Strong points of the paper}
\begin{enumerate}
    \item Numerous relevant references
    \item The introduction provides a good overview of the paper
    \item Clear practical motivation
    \item Clear and accessible language
\end{enumerate}

\section{Weak points of the paper}
\begin{enumerate}
    \item Why choose a simple MA over alternatives such as AR or exponential smoothing?
    \item The graphics are not entirely clear: why was a window size of 7 chosen?
          A concrete error metric would be helpful.
\end{enumerate}



\end{document}


